\documentclass[12pt]{article}
\usepackage{amsmath}
\usepackage{geometry} % Required for customizing margins

\geometry{margin=1in} % Sets 1-inch margins on all sides

\title{EMCH 360 Homework \#2}
\author{Jacob C. Vaught}
\date{Sep. 25, 2023}

\begin{document}

\maketitle
\section*{Problem \#1}
Bourdon gages are commonly used to measure pressure. When such a gage is attached to the closed water tank as shown in an unspecified figure, the gage reads \(5 \, \text{psi}\). The gage is positioned \(6 \, \text{in}\) above the bottom of the tank, and the top of the water in the tank is \(12 \, \text{in}\) above the gage. What is the absolute air pressure in the tank? Assume a standard atmospheric pressure of \(14.7 \, \text{psi}\).

\subsection*{Solution}
First, let's find the pressure due to the water column above the gage:

\[
P_{\text{water}} = \rho \cdot h
\]
where  
\(\rho = 62.43 \, \text{lb/ft}^3\) (density of water),  
\(h = 1 \, \text{ft} = 12 \, \text{in}\) (height of the water column above the gage).

\[
P_{\text{water}} = 62.43 \, \text{lb/ft}^3 \times 1 \, \text{ft} = 62.43 \, \text{lb/ft}^2
\]
To convert this to psi, we divide by \(144 \, \text{in}^2/\text{ft}^2\):
\[
P_{\text{water}} = \frac{62.43 \, \text{lb/ft}^2}{144 \, \text{in}^2/\text{ft}^2} = 0.434 \, \text{psi}
\]
The gage pressure is the difference between the absolute air pressure in the tank and the pressure due to the water column:
\[
P_{\text{gage}} = P_{\text{abs}} - P_{\text{water}}
\]
Solving for the absolute pressure, we get:
\[
P_{\text{abs}} = P_{\text{gage}} + P_{\text{water}} = 5 \, \text{psi} + 0.434 \, \text{psi} = 5.434 \, \text{psi}
\]
Finally, to find the absolute pressure in the tank including the atmospheric pressure, we add the standard atmospheric pressure to our calculated value:
\[
P_{\text{abs, final}} = P_{\text{abs}} + P_{\text{atm}} = 5.434 \, \text{psi} + 14.7 \, \text{psi} = 20.134 \, \text{psi}
\]

In summary, the absolute air pressure in the tank would be \(20.134 \, \text{psi}\).


\newpage
\section*{Problem \#2}
Water, oil, and an unknown fluid are contained in connected vertical tubes. Water is present at heights of \(1 \, \text{ft}\) and \(2 \, \text{ft}\) on the \(1 \, \text{inch}\) and \(2 \, \text{inch}\) sides, respectively. Above the water on the \(2 \, \text{inch}\) side is an unknown fluid with a height of \(1 \, \text{ft}\). On the \(1 \, \text{inch}\) side, above the water is oil with specific gravity \(0.9\) and height \(2 \, \text{ft}\). Determine the density of the unknown fluid.

\subsection*{Solution}
To find the density of the unknown fluid, we need to set the pressure at the base of both tubes equal to each other because the tubes are connected.

\subsubsection*{Pressure on the 1-inch Side}
The pressure at the bottom of the \(1 \, \text{inch}\) side due to water and oil is given by:
\[
P_{1\text{in}} = \rho_{\text{water}} \cdot h_{\text{water}} + \rho_{\text{oil}} \cdot h_{\text{oil}}
\]
where \(\rho_{\text{water}} = 62.43 \, \text{lb/ft}^3\), \text{ft/s}^2\) , \(h_{\text{water}} = 1 \, \text{ft}\), \(h_{\text{oil}} = 2 \, \text{ft}\), and \(\rho_{\text{oil}} = 0.9 \times 62.43 \, \text{lb/ft}^3\).
\[
P_{1\text{in}} = 62.43 \times 1 + 0.9 \times 62.43 \times 2
\]
\[
P_{1\text{in}} = 174.804 \, \text{lb/ft}^2
\]

\subsubsection*{Pressure on the 2-inch Side}
For the \(2 \, \text{inch}\) side, the pressure due to water and the unknown fluid is:
\[
P_{2\text{in}} = \rho_{\text{water}} \cdot g \cdot h_{\text{water}} + \rho_{\text{unknown}} \cdot g \cdot h_{\text{unknown}}
\]
where \(h_{\text{unknown}} = 1 \, \text{ft}\) and \(h_{\text{water}} = 2 \, \text{ft}\).

\subsubsection*{Finding the Density of the Unknown Fluid}
Equating \(P_{1\text{in}}\) and \(P_{2\text{in}}\) and solving for \(\rho_{\text{unknown}}\):
\[
\rho_{\text{unknown}} = \frac{P_{1\text{in}} - \rho_{\text{water}} \cdot g \cdot h_{\text{water}}}{g \cdot h_{\text{unknown}}}
\]
\[
\rho_{\text{unknown}} = \frac{174.804 - 124.86}{1}
\]
\[
\rho_{\text{unknown}} \approx 50.137 \, \text{lb/ft}^3
\]

\newpage
\section*{Problem \#3}
For an inclined-tube manometer as shown in a figure, the pressure in pipe A is \(0.6 \, \text{psi}\). Both pipes A and B contain water, while the gage fluid in the manometer has a specific gravity of \(2.6\). Determine the pressure in pipe B corresponding to the differential reading shown.

\subsection*{Solution}
To find the pressure at point B, we need to account for the pressure differences due to the water and the gage fluid on both sides.
\newline\newline
Firstly, let's consider the pressure at the bottom of the gage fluid on the side of pipe A.
\[
P_{\text{bottom, A}} = P_A - \rho_{\text{water}} \times h_{\text{water, A}}
\]
where \( \rho_{\text{water}} = 62.43 \, \text{lb/ft}^3 \), and \( h_{\text{water, A}} = \frac{3}{12} \, \text{ft} \).

\[
P_{\text{bottom, A}} = 0.6 \, \text{psi} - 62.43 \times \frac{3}{12 \times 144}
\]
\[
P_{\text{bottom, A}} = 0.6 \, \text{psi} - 0.108 \, \text{psi} = 0.492 \, \text{psi}
\]
Next, let's find the pressure at the bottom of the gage fluid on the side of pipe B.
\[
P_{\text{bottom, B}} = P_B - \rho_{\text{water}} \times h_{\text{water, B}} - \rho_{\text{gage}} \times h_{\text{gage}}
\]
where \( \rho_{\text{gage}} = 2.6 \times 62.43 \, \text{lb/ft}^3 \) and \( h_{\text{gage}} = \frac{4}{12} \, \text{ft} \).
Since \( P_{\text{bottom, A}} = P_{\text{bottom, B}} \),
\[
P_B = P_{\text{bottom, A}} + \rho_{\text{water}} \times h_{\text{water, B}} + \rho_{\text{gage}} \times h_{\text{gage}}
\]
\[
P_B = 0.492 \, \text{psi} + 62.43 \times \frac{3}{12 \times 144} - 2.6 \times 62.43 \times \frac{4}{12 \times 144}
\]
\[
P_B = 0.492 \, \text{psi} + 0.0441 \, \text{psi} - 0.3686 \, \text{psi} = 0.224 \, \text{psi}
\]
Thus, the pressure at point B is \(0.224 \, \text{psi}\).

\newpage
\section*{Problem \#4}
An iceberg with a specific gravity of \(0.917\) floats in the ocean with a specific gravity of \(1.025\). Determine what percent of the volume of the iceberg is underwater.

\subsection*{Solution}

\[
F_b = \text{Weight of water displaced} = \rho_{\text{ocean}} \times V_{\text{underwater}} \times g
\]
Similarly, the weight of the iceberg can be expressed as:

\[
F_i = \rho_{\text{iceberg}} \times V_{\text{total}} \times g
\]
In equilibrium, the buoyant force must balance the weight of the iceberg:

\[
F_b = F_i \implies \rho_{\text{ocean}} \times V_{\text{underwater}} \times g = \rho_{\text{iceberg}} \times V_{\text{total}} \times g
\]

\subsubsection*{Calculating Volume Underwater}
We can rearrange the equation to find \(V_{\text{underwater}}\):

\[
V_{\text{underwater}} = \frac{\rho_{\text{iceberg}}}{\rho_{\text{ocean}}} \times V_{\text{total}}
\]
Substituting the specific gravities:

\[
V_{\text{underwater}} = \frac{0.917}{1.025} \times V_{\text{total}}
\]
\[
V_{\text{underwater}} = 0.8946 \times V_{\text{total}}
\]

\subsubsection*{Percentage Underwater}
To find the percentage of the iceberg that is underwater, we simply multiply this ratio by 100:

\[
\text{Percentage underwater} = 0.8946 \times 100 = 89.46\%
\]
Thus, approximately \(89.46\%\) of the iceberg's volume is underwater.

\newpage
\section*{Problem \#5}
The inverted U-tube manometer contains oil (SG = 0.9) and water as shown in the figure. The pressure differential between pipes A and B, \( (P_A - P_B) \), is \(-7 \, \text{kPa}\). Determine the differential reading \( h \).

\subsection*{Solution}
\begin{equation}
P = \rho g h
\end{equation}
where \( P \) is the pressure, \( \rho \) is the density of the fluid, \( g \) is the acceleration due to gravity, and \( h \) is the height.

\subsection*{Determine Differential Reading \( h \)}

Firstly, let's determine the density of oil.
\begin{align*}
\rho_{\text{oil}} &= \text{SG} \times \rho_{\text{water}} \\
&= 0.9 \times 1000 \, \text{kg/m}^3 \\
&= 900 \, \text{kg/m}^3
\end{align*}
Now, let's calculate the hydrostatic pressure at point A and point B.
\newline
At point A:
\begin{align*}
P_A &= \rho_{\text{water}} \times g \times 0.2 \\
&= 1000 \, \text{kg/m}^3 \times 9.81 \, \text{m/s}^2 \times 0.2 \, \text{m} \\
&= 1962 \, \text{Pa}
\end{align*}
At point B:
\begin{align*}
P_B &= \rho_{\text{water}} \times g \times 0.3 \\
&= 1000 \, \text{kg/m}^3 \times 9.81 \, \text{m/s}^2 \times 0.3 \, \text{m} \\
&= 2943 \, \text{Pa}
\end{align*}
Since \( (P_A - P_B) = -7 \, \text{kPa} \), the pressure difference is:
\begin{align*}
P_A - P_B &= -7000 \, \text{Pa} \\
1962 \, \text{Pa} - 2943 \, \text{Pa} + \rho_{\text{oil}} \times g \times h &= -7000 \, \text{Pa} \\
-981 \, \text{Pa} + 900 \, \text{kg/m}^3 \times 9.81 \, \text{m/s}^2 \times h &= -7000 \, \text{Pa} \\
h &= \frac{-7000 \, \text{Pa} + 981 \, \text{Pa}}{900 \, \text{kg/m}^3 \times 9.81 \, \text{m/s}^2} \\
&= -0.685 \, \text{m}
\end{align*}

\subsection*{Result}
The differential reading \( h \) is \( -0.685 \, \text{m} \).
\newpage
\section*{Problem \#6}
A piston with a cross-sectional area of \(0.08 \, \text{m}^2\) is located in a cylinder containing water. An open U-tube manometer is connected to the cylinder. Given \(h_1 = 70 \, \text{mm}\) and \(h = 110 \, \text{mm}\), what is the value of the applied force \(F\) on the piston? The weight of the piston is negligible.

\subsection*{Solution}
\begin{align*}
\rho_{\text{water}} &= 1000 \, \text{kg/m}^3 \\
\rho_{\text{mercury}} &= 13600 \, \text{kg/m}^3
\end{align*}
Firstly, we calculate the pressure from the height \( h_1 \) of water and \( h \) of mercury.
\newline
Pressure due to water:
\begin{align*}
P_{\text{water}} = \rho_{\text{water}} \times g \times h_1 \\
P_{\text{water}} = 1000 \, \text{kg/m}^3 \times 9.81 \, \text{m/s}^2 \times 0.07 \, \text{m} \\
P_{\text{water}} = 686.7 \, \text{Pa}
\end{align*}
Pressure due to mercury:
\begin{align*}
P_{\text{mercury}} = \rho_{\text{mercury}} \times g \times h \\
P_{\text{mercury}} = 13600 \, \text{kg/m}^3 \times 9.81 \, \text{m/s}^2 \times 0.11 \, \text{m} \\
P_{\text{mercury}} = 14683.6 \, \text{Pa}
\end{align*}
Since the pressure at the bottom of the manometer tube must be equal, we have:
\begin{align*}
P_{\text{water}} + P_{\text{piston}} &= P_{\text{mercury}} \\
686.7 \, \text{Pa} + F / 0.08 \, \text{m}^2 &= 14683.6 \, \text{Pa} \\
F &= (14683.6 \, \text{Pa} - 686.7 \, \text{Pa}) \times 0.08 \, \text{m}^2 \\
F &= 13996.9 \, \text{Pa} \times 0.08 \, \text{m}^2 \\
F &= 1119.752 \, \text{N}
\end{align*}

\newpage
\section*{Problem \#7}
A 10-m high, 4-m wide rectangular gate is hinged at the top edge at A and is restrained by a fixed ridge at B. The gate has 7 m submerged in water and the remaining 3 m is in the air. The density of water at \(20^\circ C\) is \(9810 \, \text{N/m}^3\). We need to find:

\begin{itemize}
\item[a.] The hydrostatic force (\(F_R\)) by hydrostatic pressure on the gate in kN.
\item[b.] The horizontal force (\(F_B\)) at the ridge to open the gate at point B in kN.
\end{itemize}

\subsection*{Solution}

\subsubsection*{Calculating the Hydrostatic Force (\(F_R\))}
The hydrostatic force on the submerged part of the gate is given by the formula:

\[
F_R = \frac{1}{2} \times \gamma \times h^2 \times w
\]
where \( \gamma = 9810 \, \text{N/m}^3 \) is the specific weight of water, \( h = 7 \, \text{m} \) is the height of the submerged part, and \( w = 4 \, \text{m} \) is the width.

\[
F_R = \frac{1}{2} \times 9810 \, \text{N/m}^3 \times (7 \, \text{m})^2 \times 4 \, \text{m} = 9640820 \, \text{N} = 9640.82 \, \text{kN}
\]

\subsubsection*{Calculating the Horizontal Force (\(F_B\))}
To open the gate at point B, the horizontal force (\(F_B\)) would need to counterbalance the hydrostatic force (\(F_R\)). The hydrostatic force acts at a depth of \(\frac{2h}{3} = \frac{2\times7}{3} \, \text{m} \) from the water surface.
\newline \newline
Using the principle of moments about the hinge point A, the sum of the moments must be zero for equilibrium:

\[
F_B \times (10 \, \text{m}) = F_R \times \left( \frac{14 \, \text{m}}{3} \right)
\]

\[
F_B = \frac{F_R \times \left( \frac{14 \, \text{m}}{3} \right)}{10 \, \text{m}} = \frac{9640.82 \, \text{kN} \times \left( \frac{14 \, \text{m}}{3} \right)}{10 \, \text{m}} = 4499.05 \, \text{kN}
\]
Thus, the horizontal force required to open the gate at point B is \(4499.05 \, \text{kN}\).

\newpage
\section*{Problem\#8}
A river barge with an approximately rectangular cross-section carries a load of grain. The barge is \(30 \, \text{ft}\) wide and \(95 \, \text{ft}\) long. When unloaded, its draft (depth of submergence) is \(7 \, \text{ft}\), and with the load of grain, the draft is \(10 \, \text{ft}\). Determine: \\
(a) the unloaded weight of the barge \\
(b) the weight of the grain.

\subsection*{Solution}

\subsubsection*{Calculating the Unloaded Weight}
To find the unloaded weight of the barge (\(W_{\text{barge}}\)), we will use the formula for buoyant force, \(F_b\):

\[
F_b = \rho \cdot V
\]
Where \( \rho = 62.4 \, \text{lb/ft}^3 \) is the density of water and \( V = 30 \, \text{ft} \times 95 \, \text{ft} \times 7 \, \text{ft} \) is the volume of water displaced when the barge is unloaded.

\[
V = 30 \, \text{ft} \times 95 \, \text{ft} \times 7 \, \text{ft} = 19950 \, \text{ft}^3
\]
\[
F_b = 62.4 \, \text{lb/ft}^3 \times 19950 \, \text{ft}^3 = 1244880 \, \text{lb}
\]
The unloaded weight of the barge is equal to the buoyant force, \(W_{\text{barge}} = F_b = 1244880 \, \text{lb}\).

\subsubsection*{Calculating the Weight of the Grain}
When loaded with grain, the draft changes to \(10 \, \text{ft}\), so the new volume of water displaced (\(V_{\text{new}}\)) becomes \(30 \, \text{ft} \times 95 \, \text{ft} \times 10 \, \text{ft} = 28500 \, \text{ft}^3\).

\[
F_{b, \text{new}} = 62.4 \, \text{lb/ft}^3 \times 28500 \, \text{ft}^3 = 1778400 \, \text{lb}
\]
The weight of the grain (\(W_{\text{grain}}\)) is the difference between the new and old buoyant forces:

\[
W_{\text{grain}} = F_{b, \text{new}} - F_b = 1778400 \, \text{lb} - 1244880 \, \text{lb} = 533520 \, \text{lb}
\]

\end{document}